\documentclass[11pt,letterpaper]{article}
% isms-27000-style.tex
% Suite preamble for the ISO/IEC 27000/27001/27002 ISMS engineering suite.
%
% RECONSTRUCTION PROVENANCE: the ten uploaded suite documents load a shared
% preamble with `\input{style.tex}`, but no `style.tex` was present in the
% upload batches or anywhere in this repository's history. This file is a
% documented, suite-specific reconstruction inferred from the macros, colors
% and environments the documents actually use (colors navy/teal/graytext/pale;
% tag commands \SourceTag, \EngineeringTag, \SourceNote{..}; breakable
% tcolorbox callouts; color-striped longtables; TikZ flow diagrams). It is
% NOT the recovered original style; the original visual design is unknown.
%
% pdfLaTeX only.
\usepackage[T1]{fontenc}
\usepackage[utf8]{inputenc}
\usepackage{lmodern}
\usepackage{microtype}
\usepackage[letterpaper,margin=0.9in,headheight=15pt]{geometry}
\usepackage[table]{xcolor}
\definecolor{navy}{HTML}{14355A}
\definecolor{teal}{HTML}{0A7580}
\definecolor{graytext}{HTML}{55616E}
\definecolor{pale}{HTML}{EFF5F7}
\usepackage{amsmath,amssymb}
\usepackage{array,booktabs,longtable,tabularx,multirow}
\usepackage{enumitem}
\usepackage[most]{tcolorbox}
\usepackage{tikz}
\usetikzlibrary{arrows.meta,positioning,shapes.geometric,fit,calc}
\usepackage{needspace}
\usepackage{titlesec}
\usepackage{fancyhdr}
\usepackage{lastpage}
\usepackage{hyperref}
\hypersetup{colorlinks=true,linkcolor=navy,urlcolor=teal,citecolor=navy,pdfauthor={ISMS Engineering Documentation Suite},pdfsubject={Original supporting guidance for ISO/IEC 27000, 27001 and 27002}}
\setlength{\parindent}{0pt}
\setlength{\parskip}{5pt plus 1pt}
\setlength{\emergencystretch}{2em}
\renewcommand{\arraystretch}{1.18}
\setlist[itemize]{itemsep=2pt,topsep=3pt,leftmargin=1.5em}
\setlist[enumerate]{itemsep=2pt,topsep=3pt,leftmargin=1.7em}
\titleformat{\section}{\Large\bfseries\color{navy}}{\thesection}{0.7em}{}[\vspace{-3pt}\color{teal}\titlerule]
\titleformat{\subsection}{\large\bfseries\color{navy}}{\thesubsection}{0.6em}{}
\titleformat{\subsubsection}{\normalsize\bfseries\color{teal}}{\thesubsubsection}{0.55em}{}
\pagestyle{fancy}
\fancyhf{}
\fancyhead[L]{\sffamily\footnotesize ISO/IEC 27000 series | ISMS Engineering}
\fancyhead[R]{\sffamily\footnotesize\color{graytext} Independent companion}
\fancyfoot[L]{\sffamily\scriptsize Original explanatory material; consult the licensed standards}
\fancyfoot[R]{\sffamily\scriptsize\thepage\ / \pageref{LastPage}}
\renewcommand{\headrulewidth}{0.3pt}
% Inline provenance tags used throughout the suite: \SourceTag marks
% statements grounded in the standards; \EngineeringTag marks original
% engineering interpretation. Both are zero-argument label commands.
\newcommand{\SourceTag}{\textsf{\footnotesize\bfseries\color{navy}[Standard]}~}
\newcommand{\EngineeringTag}{\textsf{\footnotesize\bfseries\color{teal}[Engineering]}~}
\newcommand{\SourceNote}[1]{\par\smallskip{\footnotesize\color{graytext}\textit{Trace: #1.}}\par}

\begin{document}
\pagenumbering{roman}
\begin{titlepage}
\vspace*{1.2in}
{\color{teal}\rule{\textwidth}{4pt}}\par\vspace{20pt}
{\Huge\bfseries\color{navy} Application Security and Secure SDLC}\par\vspace{18pt}
{\Large\color{graytext} Security engineering gates, GHAS/CodeQL patterns and evidentiary traceability}\par\vspace{1.3in}
{\large ISO/IEC 27000:2018 $\vert$ ISO/IEC 27001:2022 $\vert$ ISO/IEC 27002:2022}\par
\vfill
\textbf{Independent educational and engineering implementation reference}\par
Version 1.0 $\cdot$ October 2026\par
{\small Original explanatory material. Source standards are copyrighted; consult licensed originals for authoritative requirements.}\par
\end{titlepage}
\tableofcontents\clearpage\pagenumbering{arabic}
\section{Relevant control foundation}
Source references: ISO/IEC 27002:2022 8.4 and 8.25--8.34, supported by 5.8, 5.21, 5.23, 8.8, 8.15--8.16 and 8.24. These describe secure development, application requirements, engineering principles, secure coding, testing, outsourcing and change-related protection. Tools and thresholds below are additional implementation proposals.
\section{Secure SDLC stage gates}
\small\rowcolors{2}{pale}{white}\begin{longtable}{p{.15\linewidth}p{.42\linewidth}p{.37\linewidth}}
\rowcolor{navy} \textcolor{white}{\textbf{Stage}} & \textcolor{white}{\textbf{Security activity}} & \textcolor{white}{\textbf{Evidence and release test}}\\
\endfirsthead
\rowcolor{navy} \textcolor{white}{\textbf{Stage}} & \textcolor{white}{\textbf{Security activity}} & \textcolor{white}{\textbf{Evidence and release test}}\\
\endhead
Discovery & Classify information, identify stakeholders, third-party and privacy constraints & Approved security requirements and asset owner\\
Design & Threat model trust boundaries and misuse cases; choose architecture principles & Review record, tests linked to threats\\
Code & Apply secure coding standards, review and secret detection & Pull-request approvals, secret scan outcome\\
Build & Pin trustworthy dependencies, assess license and known exposures, attest artifacts & SBOM, dependency report, provenance assertion\\
Verification & Run language-aware SAST, DAST and behavior/security testing as relevant & Reviewed findings with reproduction or justified disposition\\
Release & Authorize deployment, restrict credentials, record change and rollback readiness & Approvals, signed artifact and environment check\\
Operations & Observe attack signals, reassess vulnerabilities and test incident handling & Monitoring signals, remediation tickets and exercise reports\\
\end{longtable}\normalsize

\section{Recommended GitHub security design}
\begin{enumerate}
\item Use GitHub repository/organization rules to require peer review, status checks and protected branch controls, subject to licensed feature availability.
\item Configure CodeQL for supported languages using appropriate build mode; scan pull requests and default branch, and validate coverage and query configuration.
\item Apply secret scanning, push protection where supported, and a documented workflow for verified, non-sensitive or revoked credentials.
\item Assess third-party packages with dependency analysis, SBOM generation and controlled dependency update policies.
\item Treat scanner output as a triage input; record exploitability and context, false positives, approved exceptions and retest evidence.
\item Use GitHub security alerts, REST exports and controlled evidence retention to support 27001 9.1/9.2 assessment; product telemetry does not itself certify conformity.
\end{enumerate}

\section{CodeQL-style finding lifecycle}
\begin{verbatim}
DETECTED -> TRIAGED -> ASSIGNED -> REMEDIATION_PROPOSED
         -> FIX_VALIDATED -> MERGED -> DEFAULT_BRANCH_VERIFIED
         -> CLOSED
   OR -> EXCEPTION_APPROVED -> EXPIRY_REVIEW
   OR -> VERIFIED_FALSE_POSITIVE -> QUALITY_FEEDBACK

Fields: repository, branch, query, CWE, severity, affected path,
       dataflow evidence, exploitable conditions, owner, SLA,
       suppressions rationale, retest, acceptance decision.
\end{verbatim}
\section{Example sample test: injection risk}
Scenario: an API concatenates an untrusted request parameter into a SQL query. Engineering design requires parameterized queries, rejection of unauthorized dataset references and safe error handling. Positive test checks business behavior; negative test uses harmless metacharacters and verifies that the query structure is unchanged. Link to application security requirements (8.26), secure coding (8.28) and security testing (8.29).
\section{Example sample test: credential exposure}
Scenario: a token string appears in a commit. Triage verifies token type and validity without unnecessarily exposing the secret; response revokes it when exposed and assesses scope of misuse. Prevention tests check that a safely generated dummy credential pattern is blocked or alerted, and that exceptions cannot silently bypass approval. Link to 5.17, 8.4, 8.24, 8.25 and 8.29.
\section{Recommended evidence schema}
\small\rowcolors{2}{pale}{white}\begin{longtable}{p{.33\linewidth}p{.61\linewidth}}
\rowcolor{navy} \textcolor{white}{\textbf{Field}} & \textcolor{white}{\textbf{Purpose}}\\
\endfirsthead
\rowcolor{navy} \textcolor{white}{\textbf{Field}} & \textcolor{white}{\textbf{Purpose}}\\
\endhead
control\_id / risk\_id & Bidirectional traceability from requirement and risk\\
repo / commit / artifact & Immutable software and release scope\\
scan\_type / engine / configuration & Reproducible coverage and known limitations\\
run\_id / timestamp / branch & Timing and what was actually tested\\
finding / disposition / approval & Structured adjudication rather than raw counts\\
verification / remediation reference & Proof the fix was evaluated\\
evidence\_retention / access & Controlled availability to authorized reviewers\\
\end{longtable}\normalsize

\section{Security release policy (example)}
\begin{itemize}
\item Block release when required branch checks fail; document break-glass actions with named approval and expiry.
\item Require deterministic SAST/build steps for relevant components; record explicit exclusions for unsupported code.
\item Severity thresholds are internal risk decisions, not fixed ISO numerical cutoffs.
\item Confirm fixes on the relevant branch or artifact before promotion; assess default-branch state after merge.
\item Measure mean time to validate fixes and reopen recurrence, not just number of scans.
\end{itemize}

\clearpage\section*{Source and interpretation policy}
\addcontentsline{toc}{section}{Source and interpretation policy}
\textbf{Source editions:} ISO/IEC 27000:2018 (overview and vocabulary); ISO/IEC 27001:2022 (requirements and normative Annex A); ISO/IEC 27002:2022 corrected 2022-03 (control guidance). Specific clause identifiers are references to the supplied editions, not reproduced control text.\par
\textbf{Interpretation:} ``Source-grounded paraphrase'' describes a concept in the cited source. ``Engineering recommendation'' is a practical model, technology choice, owner, evidence example, frequency, template or scoring convention authored for this suite; it is not a claim that ISO prescribes that exact approach.\par
\textbf{Conformity:} Clauses 4--10 of ISO/IEC 27001 govern ISMS conformity. Applicability of controls is determined through risk treatment, with Annex A used as a completeness cross-check and exclusions justified in the Statement of Applicability (SoA). Source documents alone neither certify an ISMS nor prove control operating effectiveness.\par
\textbf{Source licensing:} Do not distribute or embed the licensed standard PDFs in repositories or distribute copied normative text. Review the editions, amendments, and certification scheme applicable to each assessment before relying on this work.\par
\end{document}
